
%---------------------------------------------------------------------
% Appendix D
%---------------------------------------------------------------------

\chapter*{Appendix D: ECSA Graduate Attributes}
\addcontentsline{toc}{chapter}{Appendix D: ECSA Graduate Attributes}
%---------------------------------------------------------------------

\section*{Project mapping to ECSA Graduate Attributes}
\phantomsection
\addcontentsline{toc}{section}{Project mapping to ECSA Graduate Attributes}

The table below provides a structured reflection on where each of the required \gls{ga} as defined by \gls{ecsa} associated with the final year project is met. This includes reference to the specific Chapter/Section (where applicable). 

\vspace{1em}


{\small
\renewcommand{\arraystretch}{1.5}
\begin{longtable}{>{\raggedright\arraybackslash}p{0.05\textwidth} >{\raggedright\arraybackslash}p{0.25\textwidth} >{\noindent\justifying\arraybackslash}p{0.6\textwidth}}
\toprule
\textbf{GA} & \textbf{Graduate Attribute Description} & \textbf{Student reflection: How this outcome was achieved during the project} \\
\midrule
\endhead

GA1 & \textit{Problem solving} &
The lack of in situ data on the underside of sea ice in the Marginal Ice Zone is stated independently of any
solution in \autoref{sec:ch1.problemStatement}. It is broken into achievable objectives in
\autoref{sec:ch1.projectobjectives}, with the scope, limitations and assumptions given in \autoref{sec:ch1.scope}.
The problem was divided into sub-problems that were solved in turn: underwater stereo calibration
(\autoref{ch:calibration}), followed by depth estimation, pose estimation and frame fusion
(\autoref{ch:reconstruction}). Each was approached from the first principles of camera geometry, refraction and
stereo matching developed in \autoref{ch:theory}, and the conclusions in
\autoref{ch:conclusions} are supported by measured results. \\
\cmidrule(lr){1-3}
GA4 & \textit{Investigations, experiments and data analysis} &
The literature review (\autoref{ch:review}) examines existing under-ice measurement and underwater three-dimensional
reconstruction methods, and its critique (\autoref{sec:ch2.critique}) motivates the stereo approach taken. The test
rig, 3D-printed targets, foam ice targets and recorded datasets are described in \autoref{sec:testplatform} and \autoref{sec:ch4.targets}. Calibration
experiments on image channel selection, detector tuning and underwater compensation are given in
\autoref{ch:calibration}. Depth, pose and fusion methods were each implemented and compared in
\autoref{ch:reconstruction}. Performance was quantified with reprojection error, absolute trajectory error against
AprilTag poses, and point-to-\acrshort{cad} error and coverage on objects of known geometry, which also verify that the code
produces accurate results. All code is documented on GitHub (\autoref{sec:appA.first}). \\
\cmidrule(lr){1-3}
GA5 & \textit{Use of engineering tools} &
The hardware (BlueROV2 Heavy, \acrshort{rpi} and stereo cameras) and the software tools are described in
\autoref{sec:ch4.hardware} and \autoref{sec:ch4.software}. Python was used for all pipeline code, with OpenCV for
tag detection, calibration and \acrshort{sgbm}, Open3D for point clouds, \acrshort{icp} and \acrshort{tsdf} fusion, COLMAP for structure from motion
and multi-view stereo, and PyTorch to run the FoundationStereo and AquaStereo models. Kalibr was used to generate
the calibration board, Onshape for \acrshort{cad}, BlueOS and Docker to operate the \acrshort{rov}, Git and GitHub for version control,
and \LaTeX{} and Zotero for the report. The limitations of each tool were identified through testing, for example
the speckle noise of \acrshort{sgbm}, the reliance of COLMAP on texture, and the smoothing of partly hidden curves by
FoundationStereo (\autoref{subsec:ch5.fusion_discussion}). \\
\cmidrule(lr){1-3}
GA6 & \textit{Professional and technical communication} &
The report follows a logical structure, from the problem and literature through theory, system design, datasets
and experiments to discussion, conclusions and recommendations (\autoref{sec:ch1.planofdevelopment}). All
equations are typeset and numbered, all figures and tables are captioned and referred to in the text, and
acronyms are defined on first use and listed in the glossary. Sources are cited in \acrshort{ieee} style, and cross-references
are hyperlinked to help the reader navigate the document. \\
\cmidrule(lr){1-3}
GA8 & \textit{Individual and collaborative teamwork} &
The project was completed individually, with the normal guidance of the supervisor, R.A. Verrinder, whose initial
project proposal was used as a reference for the problem statement. The project received \acrshort{ebe} ethics approval
(STU-EBE-2026-PSQ002570), given in Appendix~B. All assistance, including the use of
\acrshort{ai} tools, is acknowledged in the Acknowledgements. \\
\cmidrule(lr){1-3}
GA9 & \textit{Independent learning ability} &
The topic is not covered in the undergraduate curriculum, so the required knowledge was learnt independently.
This included camera modelling and calibration, refraction at flat ports, stereo matching, learnt stereo models,
structure from motion, \acrshort{icp} and \acrshort{tsdf} fusion, as well as tools such as BlueOS, Open3D, COLMAP and PyTorch. Progress was
managed against a weekly timeline, with weekly supervisor meetings and goals. The plan was adapted after examining the data
, for example by adding learnt stereo matchers and markerless pose methods were failing. \\
\bottomrule
\end{longtable}
}





